A día de hoy se han hecho progresos significativos en el campo de las redes neuronales artificiales, siendo el principal tema de actualidad en 2026 tenemos un mercado competitivo en el ámbito de los modelos de lenguaje, con empresas como OpenAI, Google, Meta, Microsoft, entre otras, compitiendo por desarrollar modelos cada vez más avanzados y eficientes. Estos modelos se utilizan en una amplia variedad de aplicaciones, desde chatbots y asistentes virtuales hasta sistemas de recomendación y análisis de sentimientos. Además, se han realizado avances en la comprensión y mitigación de los sesgos en los modelos de lenguaje, así como en la mejora de su capacidad para generar texto coherente y relevante. En resumen, el estado del arte en redes neuronales artificiales es un campo dinámico y en constante evolución, con un enfoque particular en el desarrollo de modelos de lenguaje cada vez más sofisticados y aplicables a una amplia gama de tareas. Tanto es el desarrollo actual que presenciamos como cambian incluso sectores laborales enteros, poniendo como ejemplo la salida más comun de mis estudios, el desarrollador de software, actualmente pasa un momento delicado pues cada vez más empresas optan por utilizar modelos de lenguaje para generar código, lo que ha llevado a una disminución en la demanda de desarrolladores humanos. Sin embargo, también se están creando nuevas oportunidades laborales en áreas como la supervisión y mejora de los modelos de lenguaje, así como en la investigación y desarrollo de nuevas aplicaciones para esta tecnología. En general, el estado del arte en redes neuronales artificiales es un campo emocionante y prometedor, con un impacto significativo en la forma en que interactuamos con la tecnología y el mundo que nos rodea.